\section{Auction Model and Target Parameter}

Consider auctions indexed by $j=1,\ldots,n$. In the baseline model each auction
has a known and fixed number $N\geq2$ of risk-neutral bidders. Bidder
valuations satisfy
\begin{equation}
V_{ij}\overset{\mathrm{iid}}{\sim}F_0,
\qquad i=1,\ldots,N,
\end{equation}
independently across auctions. The distribution $F_0$ is continuously
differentiable on the positive real line with density $f_0$. The observed data
are $O_j=(B_j,N)$, where
\begin{equation}
B_j=\max_{1\leq i\leq N}V_{ij}.
\end{equation}
Under truthful bidding in the second-price benchmark, $B_j$ is the highest
submitted bid. It is not the transaction price, which is the second-highest bid
when the reserve does not bind. This distinction is an empirical requirement,
not a relabeling of the observed price.

For a regular valuation distribution, the Myerson reserve $r_0$ is the unique
solution to
\begin{equation}
1-F_0(r_0)-r_0 f_0(r_0)=0.
\end{equation}
Equivalently, $r_0$ maximizes the single-bidder posted-price revenue
$r\{1-F_0(r)\}$. For $N$ bidders, expected revenue in a second-price auction
with common reserve $r$ is
\begin{align}
\mathcal R_N(r;F_0)
={}&r\{1-F_0(r)^N\}\\
&+\int_r^\infty
\left[1-F_0(t)^N-N\{1-F_0(t)\}F_0(t)^{N-1}\right]\,dt.
\end{align}
The baseline target is the scalar $r_0$. Later sections introduce smoothing
only where it is needed to construct an orthogonal estimating equation and
valid inference.
